\chapter{Related Work and Background}

The fuel-aware route optimization problem is a classic instance of the
Resource-Constrained Shortest Path (RCSP) problem with spatially varying costs, specifically
known as the \textit{Gas Station Problem} (GSP).
Below we summarize the two most directly relevant works, what each achieves, and how this
project builds upon them.

\section{Khuller, Malekian, and Mestre (2008) -- Foundational Optimality Results}

Khuller et al.~\cite{khuller2008} provide the first rigorous algorithmic treatment of the GSP.
Their central contribution is a polynomial-time algorithm with complexity
$O(\Delta n^{2} \log n)$ (where $\Delta$ is the maximum node degree) that finds the globally
optimal refueling strategy when fuel can be purchased in any amount.
Their proof of correctness relies on dynamic programming over the set of stations reachable
with a full tank, and they formally establish that the optimal substructure property holds for
this problem.

They prove that partial refueling, buying only as much fuel as needed to reach the next
cheapest station, can strictly reduce cost compared to always filling to capacity.
They quantify this gap and show the full-tank-only policy can be arbitrarily suboptimal.

Their work is primarily theoretical and does not evaluate performance on real road network
data.
Their algorithm assumes a simplified network model and does not address heuristic search
acceleration.

We implement and benchmark their DP formulation as one of our six strategies.
We directly test their theoretical prediction, that partial refueling saves 5--15\% cost over
full-tank-only, on real Thailand provincial data, providing what we believe is the first
empirical validation at this geographic scale.

\section{Nandy et al.\ (2023) -- Refuel A* (RF-A*)}

Nandy et al.~\cite{nandy2023} introduce Refuel A* (RF-A*), a heuristic search algorithm
that extends A* specifically for path finding under fuel constraints.
RF-A* iteratively constructs partial paths: at each step it identifies the next refuelling
decision point and applies A* to plan a least-cost segment to that point, respecting the
remaining fuel.
The heuristic used is:
\[
    h(v) = \frac{d_{\mathrm{geo}}(v, t)}{\eta} \cdot p_{\min},
\]
where $d_{\mathrm{geo}}$ is the straight-line (geodesic) distance to the destination and
$p_{\min}$ is the globally cheapest fuel price in the network.
This heuristic is admissible as it never overestimates the remaining cost, so RF-A* is
guaranteed to find an optimal solution.

RF-A* expands significantly fewer states than Dijkstra-based approaches on sparse graphs
because the heuristic prunes states whose estimated total cost exceeds the current
best-known solution.
They demonstrate speedups on benchmark path-finding instances.

The heuristic $h(v) = d_{\mathrm{geo}}(v,t)/\eta \cdot p_{\min}$ is \emph{fuel-level agnostic}:
it estimates the cost of traveling the remaining distance at the cheapest possible price but
does not account for how much fuel the vehicle currently has.
If the vehicle already has enough fuel to reach a cheap station and refuel there, the heuristic
may still assign high priority to states that will ultimately be inexpensive, meaning it fails to
prune them early.
Additionally, Nandy et al.\ evaluate on synthetic or small-scale networks and do not test on
real national road data with authentic price variation.

We implement RF-A* as a primary baseline and propose a direct extension described in the
next section.

\section{Our Proposed Extension: Partial-Fill A* (PF-A*)}
\label{sec:pfa_overview}

The principal algorithmic contribution of this project is \textit{Partial-Fill A*} (PF-A*), a
fuel-level-aware heuristic extension of RF-A* that addresses its fuel-agnostic limitation.
We replace $h(v)$ with a heuristic $h_{\mathrm{PF}}(v, f)$ that conditions on the vehicle's
current fuel level $f$:
\[
    h_{\mathrm{PF}}(v, f)
    = \max\!\left(\frac{d_{\mathrm{geo}}(v, t)}{\eta} - f,\; 0\right)
      \cdot p_{\min}^{\mathrm{reach}}(v, f),
\]
where $p_{\min}^{\mathrm{reach}}(v, f)$ is the minimum fuel price among all stations reachable
from $v$ with $f$ liters remaining.

The key intuitions are:
\begin{enumerate}
    \item \textbf{Fuel-sufficiency check.}
          If $f \geq d_{\mathrm{geo}}(v,t)/\eta$, the vehicle already has enough fuel to reach
          the destination in the best case; the heuristic returns~0.
    \item \textbf{Tighter price estimate.}
          Instead of the global $p_{\min}$, we use the cheapest station actually accessible from
          state $(v, f)$. On Thai networks with regional price variation,
          $p_{\min}^{\mathrm{reach}}$ is often 5--10\,THB/L higher than the global minimum,
          making $h_{\mathrm{PF}}$ a tighter lower bound.
\end{enumerate}

It is important to frame PF-A* accurately: it is a practically motivated \emph{heuristic
improvement within the existing RF-A* framework}, not a fundamentally new problem
formulation.
Its theoretical worst-case complexity is identical to RF-A*, and its optimality guarantee
follows from inherited admissibility and consistency properties.
The contribution lies in the empirical reduction of expanded states on Thailand-scale networks
and in the formal admissibility proof for the fuel-conditioned heuristic, which is developed in
Section~\ref{sec:pfa_correctness}.
